\documentclass[12pt,a4paper,twoside]{report}

\usepackage[utf8]{inputenc}
\usepackage[T1]{fontenc}
\usepackage[spanish,es-tabla]{babel}
\usepackage{mathptmx}
\usepackage{amsmath,amssymb,amsthm}
\usepackage{graphicx}
\usepackage{geometry}
\usepackage{hyperref}
\usepackage{fancyhdr}
\usepackage{setspace}
\usepackage{enumerate}
\usepackage{caption}
\usepackage{subcaption}
\usepackage{float}
\usepackage{booktabs}
\usepackage{parskip}

\geometry{top=2.5cm,bottom=2.5cm,left=2.5cm,right=2.5cm}

\onehalfspacing

\hypersetup{
    colorlinks=true,
    linkcolor=black,
    filecolor=black,
    urlcolor=blue,
    citecolor=black
}

\captionsetup{font=small,labelfont=bf}

\newtheorem{teorema}{Teorema}
\newtheorem{proposicion}{Proposición}
\newtheorem{definicion}{Definición}
\newtheorem{ejemplo}{Ejemplo}
\newtheorem{lema}{Lema}

\title{Representación de Datos Complejos en Sistemas de Información Académica}
\author{TFG - Universidad de Castilla-La Mancha}
\date{\today}

\pagestyle{fancy}
\fancyhf{}
\fancyhead[LO,LE]{\leftmark}
\fancyhead[RO,RE]{\thepage}

\begin{document}

\maketitle

\tableofcontents
\newpage

\chapter{Introducción}
\label{chap:introduccion}

\section{Contexto y motivación}
\label{sec:contexto}


\section{Objetivos}
\label{sec:objetivos}


\section{Estructura del documento}
\label{sec:estructura}


\chapter{Estado del Arte}
\label{chap:estado_del_arte}

\section{Formatos existentes para CV académicos}
\label{sec:formatos_existentes}


\section{Limitaciones y necesidades actuales}
\label{sec:limitaciones}


\chapter{Fundamentos Teóricos}
\label{chap:fundamentos}

\section{Esquemas de datos y representación de información}
\label{sec:esquemas_datos}


\section{Validación de datos con Pydantic}
\label{sec:pydantic}


\section{Normativas y estándares académicos}
\label{sec:estandares}


\chapter{Metodología}
\label{chap:metodologia}

\section{Diseño del esquema de datos}
\label{sec:diseno_esquema}


\section{Implementación de validadores}
\label{sec:implementacion_validadores}


\section{Desarrollo de herramientas de exportación}
\label{sec:herramientas_exportacion}


\chapter{Resultados}
\label{chap:resultados}

\section{Esquema propuesto}
\label{sec:esquema_propuesto}


\section{Herramientas desarrolladas}
\label{sec:herramientas_desarrolladas}


\section{Validación y pruebas}
\label{sec:validacion}


\chapter{Discusión}
\label{chap:discusion}

\section{Análisis de resultados}
\label{sec:analisis_resultados}


\section{Comparación con soluciones existentes}
\label{sec:comparacion}


\section{Limitaciones y trabajo futuro}
\label{sec:futuro}


\chapter{Conclusiones}
\label{chap:conclusiones}


\chapter{Bibliografía}
\label{chap:bibliografia}

\bibliographystyle{apalike}
\bibliography{referencias}

\end{document}
